\documentclass[11pt]{article}
\usepackage[a4paper,margin=1in]{geometry}
\usepackage{amsmath,amssymb,amsthm,mathtools}
\usepackage[hidelinks]{hyperref}
\usepackage{booktabs}

\newtheorem{theorem}{Theorem}
\newtheorem{proposition}[theorem]{Proposition}
\newtheorem{lemma}[theorem]{Lemma}
\newcommand{\Cl}{\operatorname{Cl}}
\newcommand{\Tr}{\operatorname{Tr}}
\newcommand{\Q}{\mathbb Q}
\newcommand{\Z}{\mathbb Z}

\title{A Dimension-Five Twisted-Convolution Identity\\
from a Shintani Ray Packet}
\author{Draft for review}
\date{27 July 2026}

\begin{document}
\maketitle

\begin{abstract}
For $K=\Q(\sqrt3)$ and modulus $(5)\infty_2$, the ray group relevant
to the dimension-five principal ghost is cyclic of order eight.  We
compute all twenty-four nonzero characteristics, their positive Kopp
lifts, ray classes, cocycle signs, and modular multipliers.  We then
prove unconditionally that the four positive double-sine values form
the reciprocal degree-eight ray-unit packet previously recognized
numerically by Kopp.  The proof specializes Shintani's 1978
algebraicity theorem, extracts the safe exponent $5760$, certifies the
labeled ray field and Frobenius orbit exactly, and uses rigorous Arb
enclosures together with Voutier's height bound to remove the power.
An exact calculation over $\Q(\zeta_5,\sqrt6)$ then proves that all one
hundred $2$-minors of the reconstructed ghost matrix vanish.  This
establishes both predicted shifts for every dimension-five admissible
tuple, and hence the complete dimension-five Twisted Convolution
Conjecture.  No Stark conjecture or \texttt{bnrstark} recognition is
used as an input.
\end{abstract}

\section{Canonical data and conventions}

Put
\[
 K=\Q(\sqrt3),\qquad \beta=2+\sqrt3,\qquad
 Q=\langle1,-4,1\rangle,\qquad
 L=\begin{pmatrix}4&-1\\1&0\end{pmatrix}.
\]
Then
\[
 B=L^3=\begin{pmatrix}56&-15\\15&-4\end{pmatrix},\qquad
 B\binom0{1/5}-\binom0{1/5}=\binom{-3}{-1}.
\]
Label the real embeddings by
$\infty_1(\sqrt3)=+\sqrt3$ and
$\infty_2(\sqrt3)=-\sqrt3$.  The modulus throughout is
\[
 \mathfrak m=(5)\infty_2.
\]

To remove the standard double-sine ambiguity, define
\[
 S_2(z\mid\omega_1,\omega_2)
 :=\frac{\Gamma_2(z\mid\omega_1,\omega_2)}
 {\Gamma_2(\omega_1+\omega_2-z\mid\omega_1,\omega_2)}.
\tag{1}
\]
Thus $S_2$ is the reciprocal of Kopp's
$\operatorname{Sin}_{2,\mathrm{Kopp}}$ convention.
For $0\le c\le5$ and $0\le b<5$, write
\[
 {\cal D}_{c,b}:=
 S_2\!\left(\frac{c+b\beta}{5}\,\middle|\,\beta,1\right).
\]
For $(p,q)\ne(0,0)$, let $r\equiv-p-q\pmod5$.  The principal-form
specialization of the AFK three-factor formula is
\[
 \widetilde\nu_{p,q}
 =(-1)^{5(p+q)+pq+\min(5,p+q)}
 {\cal D}_{5-p,q}{\cal D}_{5-r,p}{\cal D}_{5-q,r}.
\tag{2}
\]
Formula (2) is used only for nonzero characteristics.  Later we set the
table entry $T_{0,0}:=\sqrt6$ solely so that the identity Weyl
coefficient is $T_{0,0}/\sqrt6=1$.  This reconstruction convention is
not AFK's auxiliary zero-characteristic cocycle value.

\begin{lemma}[Derivation and convention check for (2)]
Formula \textup{(2)} is the specialization of AFK Definition~1.32 to
$t=(5,1,\langle1,-4,1\rangle)$, with no unrecorded reciprocal or sign.
\end{lemma}

\begin{proof}
AFK Definition~1.32 first forms
$\widetilde\nu_{\boldsymbol p}
=\Phi_{\boldsymbol p}\operatorname{shin}^{\boldsymbol p/5}_{B}(\beta)$.
Insert the defining modular-to-Jacobi quotient
\cite[equation (1.26)]{AFK}, and then use the Jacobi cocycle
decomposition \cite[equations (8.2)--(8.5)]{AFK} and the $S$-kernel
formula \cite[equation (8.6)]{AFK}.  In the notation of AFK Section~8.2,
\[
 L=T^4S,\qquad B=L^3,\qquad L\beta=\beta,\qquad
 j_{L^k}(\beta)=\beta^k .
\]
Thus the word $B=(T^4S)^3$ contributes exactly three $S$-kernels.
Modulo five the characteristic orbit is
\[
 (p,q)\longmapsto(r,p)\longmapsto(q,r)\longmapsto(p,q),
 \qquad r\equiv-p-q\pmod5,
\]
and the three unreduced linear arguments are
\[
 \frac{q\beta-p}{5},\qquad
 \frac{p\beta-r}{5},\qquad
 \frac{r\beta-q}{5}.
\]
Reducing their first coordinates to the representatives
$5-p,5-r,5-q$ gives respectively
$\mathcal D_{5-p,q}$, $\mathcal D_{5-r,p}$, and
$\mathcal D_{5-q,r}$.  AFK's $S$-kernel is its displayed double sine
in the denominator (AFK Section~8.2, equation for
$\sigma_S$).  That double sine is Kopp's
$\operatorname{Sin}_{2,\mathrm{Kopp}}$; hence each denominator is
precisely one factor of the reciprocal convention (1).

For completeness, apply the integral-translation law
\cite[equation (8.9)]{AFK} to the three representative changes and
multiply by the phase in AFK Definition~1.30.  Direct substitution gives
\[
 \Phi_{p,q}\,(\hbox{three translation multipliers})
 =(-1)^{E(p,q)},\qquad
 E(p,q)=5(p+q)+pq+\min(5,p+q).
\]
This is a finite parity calculation:
\[
 E(p,q)\equiv
 \begin{cases}
 pq& (p+q\le5),\\
 p+q+pq+1&(p+q\ge5)
 \end{cases}\pmod2,
\]
where the two expressions agree at $p+q=5$.  In the order
$p=0,\ldots,4$ (columns $q=0,\ldots,4$), the resulting signs are
\[
\begin{pmatrix}
 *&+&+&+&+\\
 +&-&+&-&+\\
 +&+&+&+&-\\
 +&-&+&+&+\\
 +&+&-&+&-
\end{pmatrix};
\]
the asterisk is the excluded zero characteristic.  This accounts for
all twenty-four cases and proves (2) directly from the cited
definitions.
\end{proof}

\section{Ray labels and all characteristics}

\begin{lemma}[Certified ray arithmetic]\label{lem:ray}
The class and ray groups are
\[
 \Cl(K)=1,\qquad
 \Cl_{(5)}(K)\cong C_4,\qquad
 \Cl_{\mathfrak m}(K)\cong C_8,\qquad
 \Cl_{(5)\infty_1\infty_2}(K)\cong C_8\times C_2.
\tag{3}
\]
Kopp's residue/sign class $\mathfrak R$ is the element $4\in C_8$.
The fiber of
$\Cl_{(5)\infty_1\infty_2}(K)\to\Cl_{\mathfrak m}(K)$ has order two,
so the exponent in Kopp's theorem is $n=2/2=1$.
The computation is unconditional.
\end{lemma}

\begin{proof}
PARI/GP~2.15.4 gives the structures in (3), and
\texttt{bnfcertify}=1.  The unit group is
$\{\pm\beta^n:n\in\Z\}$ and $\beta^3\equiv1\pmod5$.
Hence units congruent to $-1$ modulo $5$ are
$-\beta^{3k}$.  Their value under $\infty_2$ is
$-(2-\sqrt3)^{3k}<0$, so no global unit realizes the pair
$(-1\bmod5,+\text{ at }\infty_2)$ defining Kopp's $\mathfrak R$.
It is therefore the nontrivial sign class, namely the unique element
of order two in $C_8$.
\end{proof}

\begin{proposition}[No quadratic-character shortcut]\label{prop:noquadratic}
The Kopp partial-zeta difference has zero coefficient at the unique
order-two character of $\Cl_{\mathfrak m}(K)$.  It is supported on four
characters, all of order eight.
\end{proposition}

\begin{proof}
Write $\Cl_{\mathfrak m}(K)=\langle g\rangle\cong C_8$ and
$\mathfrak R=g^4$.  For
$\chi_k(g)=\exp(2\pi i k/8)$, Fourier inversion gives
\[
\zeta_{\mathfrak m}(s,\mathfrak A)
-\zeta_{\mathfrak m}(s,\mathfrak R\mathfrak A)
=\frac18\sum_{k=0}^7
\overline{\chi_k(\mathfrak A)}
\bigl(1-\overline{\chi_k(\mathfrak R)}\bigr)L(s,\chi_k).
\tag{CF}
\]
Since $\chi_k(\mathfrak R)=(-1)^k$, the coefficient is zero for even
$k$ and two for odd $k$.  The support is therefore
$k=1,3,5,7$, whose characters all have order eight.  The unique
quadratic character is $\chi_4$; it is even on $\mathfrak R$ and its
coefficient in \textup{(CF)} is zero.  Hence the difference cannot
factor through the quadratic quotient.  This is the precise reason
the dimension-four quadratic $L$-function reduction does not recur.
Equivalently, $\mathfrak R=g^4$ lies in the kernel
$\langle g^2\rangle$ of the quotient $C_8\to C_2$, so
$\mathfrak A$ and $\mathfrak R\mathfrak A$ have the same image in the
unique quadratic quotient and their difference projects to zero.
\end{proof}

For each $(p,q)\ne(0,0)$ start with $p$ and repeatedly subtract five;
let $\widetilde p$ be the first lift for which
\[
 q(2-\sqrt3)-\widetilde p>0.
\tag{4}
\]
Let $\mathfrak q$ be the prime above $3$, and use it as the generator
of $C_8$.  PARI's raw cyclic basis $g$ gives the discrete logarithms
of the positive principal ideals $(q\beta-\widetilde p)$ in the audit
transcript, with $\mathfrak q=g^5$.  Multiplication by $5$, which is
its own inverse modulo $8$, converts those raw logs to the normalized
$\mathfrak q$-coordinates below:
\[
\begin{array}{c|ccccc}
p\backslash q&0&1&2&3&4\\ \hline
0&-&0&6&2&4\\
1&4&7&5&7&0\\
2&2&5&5&6&3\\
3&6&7&2&1&1\\
4&0&4&3&1&3
\end{array}
\tag{5}
\]
The corresponding positive lifts are
$0$ in the first row, $-4$ in row one except at $(1,4)$ where it is
$1$, and $-3,-2,-1$ in rows two, three, and four.

\begin{lemma}[All multiplier comparisons]\label{lem:multiplier}
For every nonzero $(p,q)\bmod5$, Kopp's modular multiplier equals
the square of the AFK phase at that characteristic.
\end{lemma}

\begin{proof}
Write $B=(\begin{smallmatrix}a&b\\c&d\end{smallmatrix})$ and
$\boldsymbol r=(\widetilde p/5,q/5)^{\mathsf T}$, with
$\widetilde p$ chosen by (4).  Substitution in Kopp's defining
theta-character gives the exponent
\[
\frac12\{(c-d+1)r_1+(-a+b+1)r_2-cd r_1^2
+2(a-1)d r_1r_2-(a-2)b r_2^2\}.
\tag{6}
\]
For $B=(\begin{smallmatrix}56&-15\\15&-4\end{smallmatrix})$, direct
reduction of (6) modulo one gives
\[
 \frac{p^2-4pq+q^2}{5}\pmod1
\tag{7}
\]
in all twenty-four cases.  Also
$s(-4,15)=-19/90$, so the Rademacher invariant is three and
$\psi^2(B)=i$.  Therefore the Kopp multiplier has exponent
\[
 -\frac14-\frac{p^2-4pq+q^2}{5}\pmod1,
\tag{8}
\]
which is precisely the exponent of the square of the AFK phase.
The exact twenty-four-row substitution, including every positive lift,
is part of the bridge certificate.
\end{proof}

For a normalized ray log $k$, name the positive square root of its
ray value by
\[
\begin{array}{c|cccccccc}
k&0&1&2&3&4&5&6&7\\ \hline
\text{positive root}&x&w&y^{-1}&z&x^{-1}&w^{-1}&y&z^{-1}.
\end{array}
\]
The class log in (5) selects the unsigned entry and the explicit sign
in (2) selects its sign.  Thus the complete reconstruction table is
\[
T_5=
\begin{pmatrix}
\sqrt6&x&y&y^{-1}&x^{-1}\\
x^{-1}&-z^{-1}&w^{-1}&-z^{-1}&x\\
y^{-1}&w^{-1}&w^{-1}&y&-z\\
y&-z^{-1}&y^{-1}&w&w\\
x&x^{-1}&-z&w&-z
\end{pmatrix}.
\tag{9}
\]
Here $x,y,z,w$ denote the four positive representatives of
the ray-class pairs
\[
(0,4),\qquad(6,2),\qquad(3,7),\qquad(1,5),
\tag{10}
\]
respectively.  In particular, the table is generated from the ray logs
and the sign exponent; it is not stored as a list of expected answers in
the bridge generator.

\section{The analytic ray packet}

Set
\[
\begin{aligned}
P(U)={}&U^8-(8+5\sqrt3)U^7+(53+30\sqrt3)U^6\\
&-(156+90\sqrt3)U^5+(225+130\sqrt3)U^4\\
&-(156+90\sqrt3)U^3+(53+30\sqrt3)U^2\\
&-(8+5\sqrt3)U+1.
\end{aligned}
\tag{11}
\]
It is the reciprocal lift $P(U)=U^4F(U+U^{-1})$ of
\[
\begin{aligned}
F(V)={}&V^4-(8+5\sqrt3)V^3+(49+30\sqrt3)V^2\\
&-(132+75\sqrt3)V+(121+70\sqrt3).
\end{aligned}
\tag{12}
\]
Kopp previously recognized the same polynomial numerically as the
Stark-unit polynomial for $\Q(\sqrt3)$ and modulus five; see
\cite[\S7, (7.21)]{Kopp2021} and
\cite[Example~1.17, (3.12)]{Kopp2023}.  PARI's \texttt{bnrstark}
reproduces (12), but that recognition is not used as a proof below.

The square roots of the reciprocal packet have absolute polynomial
\[
\begin{aligned}
{\cal Q}(X)={}&X^{32}-16X^{30}+95X^{28}-260X^{26}+355X^{24}
-348X^{22}\\
&+388X^{20}-300X^{18}+195X^{16}-300X^{14}+388X^{12}\\
&-348X^{10}+355X^8-260X^6+95X^4-16X^2+1.
\end{aligned}
\tag{13}
\]
Define $w^{-1},z^{-1},x^{-1},y^{-1},y,x,z,w$ to be the unique
positive roots of (13) in, respectively,
\[
\begin{gathered}
(0.424835,0.424836),\ (0.481476,0.481477),\
(0.506963,0.506964),\ (0.782390,0.782391),\\
(1.278133,1.278134),\ (1.972526,1.972527),\
(2.076946,2.076947),\ (2.353849,2.353850).
\end{gathered}
\tag{14}
\]
Sturm's theorem certifies exactly one root in every interval.

\subsection{Shintani's theorem with an explicit exponent}

For a ray class $A\in\Cl_{\mathfrak m}(K)$ put
\[
 X_A:=\exp Z'_{\mathfrak m}(0,A).
\]

\begin{lemma}[Specialized Shintani algebraicity]\label{lem:shintani}
For every $A\in\Cl_{\mathfrak m}(K)$,
\[
 X_A^{5760}\in\mathcal O_H^\times,
\]
where $H$ is the ray field of modulus $\mathfrak m$.  These powers have
the Artin action indexed by the ray classes.  At every embedding of
$H$ inducing $\infty_2$ on $K$, their absolute value is one.
\end{lemma}

\begin{proof}
Work first in the full narrow ray group
\[
 C:=\Cl_{(5)\infty_1\infty_2}(K)\cong C_8\times C_2.
\]
In Shintani's notation, let $\mu$ be the class represented by
$11-5\beta$, which is congruent to $1$ modulo $5$ and has signs
$(-,+)$ at $(\infty_1,\infty_2)$.  Let $\nu$ be the class represented
by $4$, which is congruent to $-1$ modulo $5$ and is totally positive.
The exact ray logs are
\[
 \mu=[4,1]\in C,\qquad \nu=[4,0]\in C.
\]
Forgetting $\infty_1$ kills $\mu$ and sends $\nu$ to the unique
order-two element $4\in C_8$.  Hence
\[
 G=\langle\mu\rangle,\qquad \nu\notin G,\qquad
 C/G\simeq\Cl_{\mathfrak m}(K).
\]

Every unit is $\pm\beta^n$, and $\beta$ has order three modulo $5$.
Thus no totally positive unit is $-1$ modulo $5$, which is Shintani's
condition (0-3).  There is no norm-$-1$ unit, since
$a^2-3b^2=-1$ is impossible modulo $3$; this implies condition (0-6).

The field $H$ has signature $(8,4)$ and a unique octic subfield
\[
 M\simeq\Q(\zeta_{60})^+.
\]
Thus $H/M$ is quadratic, $M$ is the maximal absolutely abelian
subfield of $H$, and exactly one real place of $K$ splits in $H$.
These are Shintani's hypothesis (0-9).  His Theorem~2
\cite{Shintani1978} therefore applies to $X_{(5)}(c,G)$.

This invariant is exactly the one-place zeta difference used here.
Indeed, if $c$ maps to $A$, then the inverse images of $A$ and
$RA$ in $C$ are respectively
\[
 \{c,c\mu\},\qquad \{c\nu,c\mu\nu\}.
\]
Shintani's Theorem~1 consequently gives
\[
\begin{aligned}
\log X_{(5)}(c,G)
&=\zeta'_{\rm full}(0,c)-\zeta'_{\rm full}(0,c\nu)\\
&\quad+\zeta'_{\rm full}(0,c\mu)
       -\zeta'_{\rm full}(0,c\mu\nu)\\
&=\zeta'_{\mathfrak m}(0,A)-\zeta'_{\mathfrak m}(0,RA)
=Z'_{\mathfrak m}(0,A).
\end{aligned}
\]

It remains to make Shintani's existential power explicit.  The normal
closure of $H$ is the full two-infinite-place ray field.  It is also
the degree-sixteen ray field over $\Q(\sqrt{-5})$ with conductor
\[
 \mathfrak c=\begin{pmatrix}15&0\\0&3\end{pmatrix}
\]
in PARI's integral basis.  The quadratic-subfield enumeration contains
a unique copy of $\Q(\sqrt{-5})$, so this identification preserves the
imaginary quadratic field occurring in Shintani's proof.

Shintani's Proposition~4 expresses the real invariant as a product of
ratios of imaginary-quadratic $W$-invariants with integral
multiplicities.  For a nontrivial divisor $\mathfrak d\mid\mathfrak c$,
clearing the absolute-value square root and the distribution exponent
requires a multiple of $12f_{\mathfrak d}n(S)$.  For the trivial
divisor the required exponent is $12h_kn(S)$.  The exact exponents for
the eight divisors are
\[
 384,\ 288,\ 288,\ 144,\ 240,\ 360,\ 360,\ 180,
\]
whose least common multiple is $5760$.  Since $(5)$ is prime and its
only proper divisor fails (0-3), Shintani's $Y_{(5)}$ equals
$X_{(5)}$ and no induction exponent is added.  Proposition~5 supplies
the unit, Artin, and nonsplit-modulus-one assertions.
\end{proof}

\subsection{The labeled candidate field}

The absolute polynomial of a root of (11) is
\[
\begin{aligned}
p(T)={}&T^{16}-16T^{15}+95T^{14}-260T^{13}+355T^{12}
-348T^{11}\\
&+388T^{10}-300T^9+195T^8-300T^7+388T^6-348T^5\\
&+355T^4-260T^3+95T^2-16T+1.
\end{aligned}
\]

\begin{lemma}[Exact field and Frobenius labels]\label{lem:frobenius}
The field defined by (11) is the labeled ray field $H$ over $K$.
Let $U$ be the root of $p$ in
\[
 (3.890,3.891).
\]
Then $U$ is a norm-one unit.  Arithmetic Frobenius at the prime
$\mathfrak q$ above $3$ acts transitively on its eight real conjugates,
which lie successively in
\[
\begin{gathered}
(3.890,3.891),\ (5.540,5.541),\ (0.612,0.613),\
(4.313,4.314),\\
(0.257,0.258),\ (0.180,0.181),\ (1.633,1.634),\
(0.231,0.232).
\end{gathered}
\]
At the embeddings above $\infty_2$, all conjugates of $U$ have
absolute value one.
\end{lemma}

\begin{proof}
Applying \texttt{rnfequation} to both the candidate and the ray field,
\texttt{nfisisom} returns eight maps.  Exact substitution shows that
all eight maps fix the labeled root of $y^2-4y+1$; hence the
isomorphism is over $K$, not merely over $\Q$.  Moreover
\texttt{bnfcertify}=1, \texttt{bnfisunit} gives an exact unit
decomposition, and the norm of $U$ is one.

The prime $\mathfrak q$ generates the ray group.  Modulo
$\mathfrak q$ the relative polynomial is
\[
 T^8+T^7+2T^6+2T^2+T+1\in\mathbb F_3[T],
\]
which is irreducible.  Arithmetic Frobenius is therefore $T\mapsto
T^3$.  Reduction modulo this polynomial identifies one of the exact
\texttt{nfgaloisconj} automorphisms with Frobenius.

A rational interval of width $10^{-80}$ contains exactly one root of
$p$ in $(3.890,3.891)$.  Exact rational interval evaluation of the
successive Frobenius polynomials places their images inside the eight
disjoint intervals in the statement.  This proves the labels
Galois-equivariantly.

Finally, under the nonidentity embedding of $K$, substitute
$T=2(Y^2-1)/(Y^2+1)$ in the trace polynomial (12).
The exact \texttt{nfpolsturm} counts are $[8,0]$, so all four trace
roots at that embedding lie in $(-2,2)$.  Each corresponding
reciprocal quadratic has both roots on the unit circle.
\end{proof}

\subsection{Certified removal of the power}

\begin{theorem}[Unconditional ray packet]\label{thm:packet}
For
\[
 \mathcal X_k:=\exp Z'_{\mathfrak m}(0,\mathfrak q^k),
 \qquad 0\le k<8,
\]
the exact Artin-labeled packet is
\[
(\mathcal X_0,\ldots,\mathcal X_7)
=(x^2,w^2,y^{-2},z^2,x^{-2},w^{-2},y^2,z^{-2}).
\tag{15}
\]
In particular, $\mathcal X_0=U$.
\end{theorem}

\begin{proof}
Kopp's proved Kronecker-limit theorem, with the convention and
multiplier checked above, gives
\[
 \widetilde\nu_{p,q}^{\,2}
 =\exp Z'_{\mathfrak m}(0,\mathfrak A_{p,q});
\tag{16}
\]
here $\mathfrak A_{p,q}$ is the normalized ray class in (5).
The four positive generators at characteristics
\[
 (0,1),\quad(0,2),\quad(2,4),\quad(3,3)
\]
have $\mathfrak q$-coordinates $0,6,3,1$, respectively.

We evaluate the regularized integral for the reciprocal double sine
(1) with Arb balls through \texttt{python-flint}~\cite{pythonflint}.
Every Simpson panel uses an Arb enclosure of the
fourth derivative.  All endpoints, widths, tolerances, and acceptance
comparisons are exact rationals or outward-rounded balls.  Near zero,
the positive series of
$H_c(t)=\sinh(ct/2)/(ct/2)$ gives the rational bound
\[
 |f(t)-f(0)|\le4t^2\qquad(0\le t\le10^{-4}).
\]
For the tail, every reduced argument and its complement belongs to
$(1/100,5)$; beyond $36$, each exponential denominator exceeds
$3/4$, giving the remainder bound $e^{-36}/9$.
The candidate roots are selected by disjoint rational isolating
intervals, independently of library ordering.

For each of the four characteristics, the ball for
$\log(\widetilde\nu_{p,q}^{\,2})-\log(U_k)$ contains zero.  The largest
radius is less than
\[
 \epsilon=2.23\times10^{-9}.
\]
The reciprocal partners give the remaining four real embeddings.

By Lemma~\ref{lem:shintani}, both
$\mathcal X_0^{5760}$ and $U^{5760}$ are units in the same labeled
field.  Set
\[
 \eta=\mathcal X_0^{5760}/U^{5760}.
\]
Lemma~\ref{lem:frobenius} and the interval labels show that at every
real embedding,
\[
 |\log|\sigma(\eta)||<5760\epsilon<1.29\times10^{-5}.
\]
At the four complex pairs, both numerator and denominator have
absolute value one.  Hence
\[
 h(\eta)<1.29\times10^{-5}.
\]

If $3\le d=[\Q(\eta):\Q]\le16$ and $\eta$ is not a root of unity,
Voutier's explicit bound \cite{Voutier1996} gives
\[
 h(\eta)>
 \frac1{4d}\left(\frac{\log\log d}{\log d}\right)^3.
\]
Its minimum for $3\le d\le16$ is
$5.2279533222\ldots\times10^{-5}$, a contradiction.  A non-torsion
quadratic unit has height at least
$\tfrac12\log((1+\sqrt5)/2)$, and a rational unit is $\pm1$.
Thus $\eta$ is a root of unity.  Since $H$ has a real embedding, its
only roots of unity are $\pm1$; positivity at the distinguished
embedding gives $\eta=1$.  Therefore
$\mathcal X_0^{5760}=U^{5760}$ and positivity yields
$\mathcal X_0=U$.  The exact Artin action in
Lemma~\ref{lem:shintani}, together with the certified Frobenius orbit
in Lemma~\ref{lem:frobenius}, proves the full packet (15).
\end{proof}

Since every reciprocal double-sine factor in (2) is positive at the
displayed arguments, equations (2), (5), and (15) now derive the
signed table (9) unconditionally.

\section{Exact finite cancellation}

Let $\zeta_5=e^{2\pi i/5}$ and
$\tau=-e^{\pi i/5}=\zeta_5^3$.  Define, without reference to software,
\[
 K_{r,c}:=\frac1{5\sqrt6}
 \sum_{\substack{0\le a,b<5\\r\equiv c+a\pmod5}}
 T_{a,b}\tau^{ab}\zeta_5^{bc},
 \qquad 0\le r,c<5,
\tag{17}
\]
where $T$ is (9).  Equivalently,
$K=\frac15\sum_{a,b}(T_{a,b}/\sqrt6)D_{a,b}$ for
$D_{a,b}e_c=\tau^{ab}\zeta_5^{bc}e_{c+a}$.
In particular,
\[
 \Tr K=T_{0,0}/\sqrt6=1.
\tag{18}
\]

\begin{proposition}[Exact 100-minor certificate]\label{prop:minors}
For the algebraic values selected by (13)--(14), all one hundred
$2$-minors of the matrix (17) vanish exactly.  Hence
$\operatorname{rank}K=1$ and $K^2=K$.
\end{proposition}

\begin{proof}
Over the coefficient field $\Q(\zeta_5,\sqrt6)$, (13) has four
degree-eight factors.  Let $w$ be the interval-defined root in (14).
Exact rational interval arithmetic applied to the
\texttt{nfgaloisconj} polynomials proves that
\texttt{nfgaloisconj} substitutions are
\[
\begin{array}{c|cccccccc}
 &w&w^{-1}&y&y^{-1}&x&x^{-1}&z&z^{-1}\\ \hline
\text{conjugate index}&2&8&5&4&10&6&16&12.
\end{array}
\tag{19}
\]
with every image lying inside its corresponding interval in (14).
Moreover, exact \texttt{nfisincl} identities express $\sqrt5$ and
$\sqrt6$ as polynomials in $w$; rational interval evaluation selects
their positive signs.  Among the four factors, factor four is the
unique factor on which both of these subfield elements equal the
positive $\sqrt5$ and positive $\sqrt6$.  This certifies, rather than
numerically assumes, that the interval-defined tuple is the tuple used
in the factor-four calculation.

Substitution in (9) and (17), followed by exact reduction modulo the
selected factor, makes all one hundred minors zero.  As a factor
selection check, factors one, two, and three leave respectively
$70,100,100$ nonzero minors.  This is exact number-field arithmetic,
not floating-point evaluation.  The complete script and transcript
are supplied with the manuscript.

The matrix is nonzero because (18) gives trace one.  Thus its rank is
exactly one.  Any rank-one matrix satisfies
$K^2=(\Tr K)K$, proving idempotency.
\end{proof}

\section{Dimension-five twisted convolution}

\begin{lemma}[Weyl coefficient bridge]\label{lem:bridge}
The matrix $K$ of (17) is the standard-basis Weyl reconstruction of
the canonical AFK ghost for
$t=(5,1,\langle1,-4,1\rangle)$.  Moreover, $K^2=K$ is equivalent to
the $\lambda=1$ instance of AFK equation~(1.49):
\[
\sum_{\boldsymbol q\in{\cal I}_{\boldsymbol p}}
\zeta_5^{\langle\boldsymbol p,(I+L_{z,t})\boldsymbol q\rangle}
\operatorname{shin}_{A_t}^{\boldsymbol q/5}(\beta)
\operatorname{shin}_{A_t^{-1}}^{(\boldsymbol q-\boldsymbol p)/5}(\beta)
=25\delta_{\boldsymbol p,\boldsymbol0}^{(5)}.
\tag{20}
\]
\end{lemma}

\begin{proof}
For $\boldsymbol p\ne\boldsymbol0$, put
$\nu_{\boldsymbol p}:=\widetilde\nu_{\boldsymbol p}$ as given by
(2).  The first two unconditional parts of AFK's construction give
\[
 \nu_{\boldsymbol p}\nu_{-\boldsymbol p}=1.
\tag{21}
\]
Theorem~\ref{thm:packet}, followed by (2) and (5), identifies these
twenty-four numbers with the nonzero entries of $T$.  The actual ghost
reconstruction is therefore
\[
 K=\frac15I+\frac1{5\sqrt6}
 \sum_{\boldsymbol p\ne\boldsymbol0}
 \nu_{\boldsymbol p}D_{\boldsymbol p}.
\tag{22}
\]
The artificial entry $T_{0,0}=\sqrt6$ packages the first term of
(22) into the single entry formula (17); it is not an analytic
zero-characteristic value.

Now expand $K^2-K$ using
$D_{\boldsymbol p}D_{\boldsymbol q}
=\tau^{\langle\boldsymbol p,\boldsymbol q\rangle}
D_{\boldsymbol p+\boldsymbol q}$ and compare every Weyl coefficient.
For $\boldsymbol p\ne\boldsymbol0$, multiplication by $150$ makes its
coefficient
\[
 \sum_{\substack{\boldsymbol q\bmod5\\
 \boldsymbol q\ne\boldsymbol0,\boldsymbol p}}
 \tau^{\langle\boldsymbol p,\boldsymbol q\rangle}
 \nu_{\boldsymbol p-\boldsymbol q}\nu_{\boldsymbol q}
-3\sqrt6\,\nu_{\boldsymbol p}.
\tag{23}
\]
Using (21), the product in (23) is
$\nu_{\boldsymbol q}/\nu_{\boldsymbol q-\boldsymbol p}$.

AFK's $\nu_{\boldsymbol0}^{\rm aux}$ is an auxiliary cocycle value,
not the identity coefficient in (22).  Their zero-characteristic
lemma specializes to
\[
 \nu_{\boldsymbol0}^{\rm aux}
+(\nu_{\boldsymbol0}^{\rm aux})^{-1}
=-(5-2)\sqrt6=-3\sqrt6.
\tag{24}
\]
After adjoining the two omitted endpoints
$\boldsymbol q=\boldsymbol0,\boldsymbol p$, their combined contribution
is
\[
 \nu_{\boldsymbol p}\!
 \left(\nu_{\boldsymbol0}^{\rm aux}
+(\nu_{\boldsymbol0}^{\rm aux})^{-1}\right)
=-3\sqrt6\,\nu_{\boldsymbol p}.
\]
Thus (23) is exactly the full endpoint-corrected sum
\[
 \sum_{\boldsymbol q\in{\cal I}_{\boldsymbol p}}
 \tau^{\langle\boldsymbol p,\boldsymbol q\rangle}
 \frac{\nu_{\boldsymbol q}}
 {\nu_{\boldsymbol q-\boldsymbol p}},
\tag{25}
\]
with the two endpoint ratios interpreted using
$\nu_{\boldsymbol0}^{\rm aux}$.

Write $\Phi_{\boldsymbol q}$ for the AFK phase.  Their phase-ratio
calculation specializes, with $r=f_j/f=1$ and $d_j=d=5$, to
\[
 \frac{\Phi_{\boldsymbol q}}
 {\Phi_{\boldsymbol q-\boldsymbol p}}
 =
 (-1)^{1+s_5(\boldsymbol p)}
 \tau^{Q(\boldsymbol p)}
 \tau^{\langle\boldsymbol p,\boldsymbol q\rangle}
 \zeta_5^{\langle\boldsymbol p,L_{z,t}\boldsymbol q\rangle}.
\tag{26}
\]
The first two factors are nonzero and independent of
$\boldsymbol q$.  Multiplying the remaining $\tau$ factor by the
displacement chirp in (25), and using $\tau^2=\zeta_5$, gives exactly
\[
 \zeta_5^{\langle\boldsymbol p,
 (I+L_{z,t})\boldsymbol q\rangle}.
\]
The cocycle inverse identity
\[
 \frac{\operatorname{shin}_{A_t}^{\boldsymbol q/5}(\beta)}
 {\operatorname{shin}_{A_t}^{(\boldsymbol q-\boldsymbol p)/5}(\beta)}
 =
 \operatorname{shin}_{A_t}^{\boldsymbol q/5}(\beta)
 \operatorname{shin}_{A_t^{-1}}^{
   (\boldsymbol q-\boldsymbol p)/5}(\beta)
\tag{27}
\]
gives the two factors printed in (20).  Here
$A_t=B$, $L_{z,t}=L$, and
$f_t(1)=2+9\equiv1\pmod5$, so the identity twist produces exactly the
chirp in (25).  The remaining factor depends on
$\boldsymbol p$ but is nonzero and independent of
$\boldsymbol q$.  Hence the vanishing of every nonzero Weyl coefficient
is equivalent to (20) for $\boldsymbol p\ne\boldsymbol0$.
For $\boldsymbol p=\boldsymbol0$, the cocycle inverse identity makes
the sum $25$ automatically.  This is the specialization of AFK
equations (5.78)--(5.98), including rather than suppressing their
endpoint correction.
\end{proof}

\begin{theorem}[Complete dimension-five TCC]\label{thm:tcc}
For every dimension-five admissible tuple $t$,
\[
 0,1\in{\cal Z}_t.
\]
Furthermore ${\cal Z}_{t'}={\cal Z}_t$ for any two dimension-five
admissible tuples $t,t'$.
\end{theorem}

\begin{proof}
Proposition~\ref{prop:minors} and Lemma~\ref{lem:bridge} prove the
$\lambda=1$ identity for the canonical tuple.  The twist condition is
\[
 \det(G)f_t(\lambda)\equiv1\pmod5,\qquad
 f_t(\lambda)=2\lambda+9.
\]
Thus $G=I$ gives $2\lambda+9\equiv1\pmod5$, or
$\lambda\equiv1\pmod5$.

Complex conjugation preserves idempotency and sends
$\lambda$ to
$\bar\lambda=-(\lambda+d-1)=1-\lambda\pmod5$.
It therefore sends the verified shift $1$ to shift $0$.
Equivalently the conjugate twist has determinant $-1$, and
$-(2\lambda+9)\equiv1\pmod5$ gives $\lambda\equiv0\pmod5$.

Finally, dimension-five admissibility forces $r=1$,
$K=\Q(\sqrt3)$, and $j=m=1$.  Since the relevant conductor is one,
every admissible form has discriminant twelve.  The ordinary class
number is one, so primitive indefinite forms of discriminant twelve
constitute one $\mathrm{GL}_2(\Z)$-equivalence class.  (The narrow class
number is not needed: determinant-$-1$ transformations are allowed.)
Thus every admissible form is integrally equivalent to the canonical
one in the sense required here.

For clarity, the transport does not invoke AFK's conditional
transformation theorem.  If $t_M$ is obtained from $t$ by
$M\in\mathrm{GL}_2(\Z)$, then unconditionally
\[
 A_{t_M}=M^{-1}A_tM,\qquad
 L_{z,t_M}=M^{-1}L_{z,t}M,\qquad
 \beta_{t_M}=M^{-1}\beta_t.
\]
Kopp's covariance formula sends each cocycle factor at
$(t_M,\boldsymbol q)$ to the corresponding factor at
$(t,\ell M\boldsymbol q)$, with
$\ell=\operatorname{sgn}(j_{M^{-1}}(\beta_t))$; when
$\det M=-1$ both factors are complex-conjugated, which does not change
the real TCC sum.  Reindexing
\[
 \boldsymbol q'=\ell M\boldsymbol q,\qquad
 {\cal I}'_{\ell M\boldsymbol p}=\ell M{\cal I}_{\boldsymbol p},
\]
and using symplectic covariance transforms the complete left side of
(20) for $(t_M,\boldsymbol p)$ into that for
$(t,\ell M\boldsymbol p)$.  Hence
${\cal Z}_t\subseteq{\cal Z}_{t_M}$.  Applying the same calculation to
$M^{-1}$ gives the reverse inclusion and therefore equality.
\end{proof}

\section{Scope and reproducibility}

Theorem~\ref{thm:tcc} proves the complete formal TCC in dimension five.
It does not prove a dimension-five SIC, the TCC in arbitrary
dimensions, or Zauner's conjecture.  The new number-theoretic step is
Theorem~\ref{thm:packet}: it proves the previously conjectural
dimension-five ray packet by combining Kopp's proved limit formula with
the applicable case of Shintani's proved algebraicity theorem and an
effective height argument.

The supplementary package contains:
\begin{itemize}
\item the 25-characteristic bridge certificate and generator;
\item the exact $C_8$ Fourier-support audit excluding the quadratic
character;
\item the PARI transcript with \texttt{bnfcertify}=1 and all positive
ray-class logs;
\item the Sturm root-isolation transcript;
\item the exact interval/conjugate/factor embedding certificate;
\item the formal 100-minor Laurent certificate; and
\item the exact four-factor minor-reduction transcript;
\item the Shintani sign-class, conductor, and exponent-$5760$ audit;
\item the labeled $K$-isomorphism, unit, Sturm, and Frobenius-orbit
certificate; and
\item the rigorous Arb double-sine enclosure and height-gap transcript.
\end{itemize}
Their SHA-256 digests are recorded in
\texttt{certificates/SHA256SUMS}.

\begin{thebibliography}{12}
\bibitem{Kopp}
G.~Kopp, \emph{The Shintani--Faddeev modular cocycle: Stark units from
$q$-Pochhammer ratios}, arXiv:2411.06763.
\bibitem{AFK}
D.~M.~Appleby, S.~T.~Flammia, and G.~S.~Kopp,
\emph{A constructive approach to Zauner's conjecture via the Stark
conjectures}, arXiv:2501.03970.
\bibitem{Kopp2018}
G.~Kopp, \emph{SIC-POVMs and the Stark conjectures},
Int.\ Math.\ Res.\ Not.\ IMRN 2021, no.~18, 13812--13838;
arXiv:1807.05877.
\bibitem{Kopp2021}
G.~S.~Kopp, \emph{Indefinite zeta functions},
Res.\ Math.\ Sci.\ \textbf{8}, 17 (2021),
\url{https://doi.org/10.1007/s40687-021-00252-9}.
\bibitem{Kopp2023}
G.~S.~Kopp, \emph{A Kronecker limit formula for indefinite zeta
functions}, Res.\ Math.\ Sci.\ \textbf{10}, 24 (2023),
\url{https://doi.org/10.1007/s40687-023-00384-0}.
\bibitem{Shintani1978}
T.~Shintani, \emph{On certain ray class invariants of real quadratic
fields}, J.\ Math.\ Soc.\ Japan \textbf{30} (1978), 139--167,
\url{https://doi.org/10.2969/jmsj/03010139}.
\bibitem{Voutier1996}
P.~Voutier, \emph{An effective lower bound for the height of algebraic
numbers}, Acta Arith.\ \textbf{74} (1996), 81--95,
\url{https://doi.org/10.4064/aa-74-1-81-95}.
\bibitem{PARI}
The PARI Group, \emph{PARI/GP version 2.15.4},
Université de Bordeaux, 2023,
\url{https://pari.math.u-bordeaux.fr/}.
\bibitem{pythonflint}
F.~Johansson and contributors, \emph{python-flint 0.9.0},
\url{https://github.com/flintlib/python-flint}.
\end{thebibliography}

\end{document}
